\documentclass[11pt,a4paper]{article}
\makeatletter
\def\input@path{{../v0.9/}}
\makeatother
\usepackage{jheppub}
\usepackage{booktabs}
\usepackage{siunitx}
\usepackage{microtype}
\hypersetup{pdftitle={Momentum-enhanced tau-neutrino propagation through the Earth and the UHE constraint},pdfauthor={Subhendu Rakshit}}

\title{Momentum-enhanced $\nu_\tau$ propagation through the Earth\\
and the ultra-high-energy constraint}
\author{Subhendu Rakshit}

\abstract{We study whether non-standard neutrino interactions can produce an
observable modification of ultra-high-energy $\nu_\tau$ propagation through
the Earth.  After validating the coupled $\nu_\tau$--$\tau$ transport machinery
against classic regeneration calculations, we first examined smooth
charge-radius and derivative interaction benchmarks and found their detector-
level distortions too small or too degenerate with astrophysical normalization
and slope.  We then turned to a genuinely momentum-growing dimension-seven
neutrino--quark operator.  A PDF-level calculation with CT10nlo and the
conservative pointwise EFT-domain restriction $\sqrt{\hat s}<\Lambda$ gives
$\sigma_{\rm NP}/\sigma_{\rm SM}=0.35$--$0.72$ at $100$--$300$ PeV for
representative $\Lambda=2.0$--$2.44$ TeV, demonstrating that momentum growth
can overcome the earlier smooth-model suppression.  The same enhancement,
however, persists into the ultra-high-energy regime constrained by IceCube.
Using the official 12.6-year IceCube EHE effective area and the public
KM3NeT bright-track response, we validate the KM3NeT Ehlert-spectrum benchmark
to better than $10^{-4}$ relative accuracy and obtain a corrected IceCube
Standard-Model fold of $17.67$ events for the unscaled benchmark spectrum.
A deliberately crude common-flux screen then drives the KM3NeT NC-equivalent
response proxy to approximately $0.031$ events almost independently of the
EFT cutoff.  Because the operator is neutral-current-like, this proxy is not a
direct muon-track prediction.  The central result of this checkpoint is
therefore not a no-go theorem: momentum enhancement succeeds at the
cross-section and propagation level, but it creates a new tension with the
UHE/cosmogenic-neutrino observations.  The public IceCube release does not
contain the reconstructed-bin response needed for an exact BSM EHE likelihood,
so the event-level conclusion remains a controlled screening result.}

\keywords{Tau neutrinos, neutrino propagation, neutrino deep-inelastic
scattering, effective field theory, IceCube, KM3NeT}

\begin{document}
\maketitle
\flushbottom
\tableofcontents

\section{Introduction and scientific goal}
High-energy tau neutrinos propagating through the Earth are qualitatively
different from a simply absorbed beam.  Charged-current production of taus,
followed by tau decay, regenerates $\nu_\tau$ at lower energy; neutral-current
scattering redistributes the spectrum; secondary neutrinos from tau decay can
feed other flavors.  The resulting energy--zenith structure is a sensitive
probe of the neutrino interaction kernel rather than merely its total
normalization.

The aim of this project is to determine whether a physically allowed
non-standard $\nu_\tau$ interaction can generate an observable Earth-
propagation distortion.  This is not a search for a no-go theorem, and it is
not the separate problem of inferring the UHECR proton fraction from a
cosmogenic-neutrino flux.  Cosmogenic/UHE neutrino data enter here because any
interaction that grows rapidly with momentum must remain consistent with the
same neutrinos at still higher energies.

The work proceeded in stages.  The transport machinery was first validated
against historical $\nu_\tau$ regeneration calculations.  Smooth
charge-radius and forward-blind derivative interactions were then propagated
through IceCube- and ARCA-like response models; their observable distortions
were too weak or too degenerate with flux normalization and slope.  That
failure motivated the present momentum-growing operator.  The momentum-
enhanced calculation succeeds in producing a large enough high-energy cross-
section modification.  The new issue is therefore not whether the interaction
can grow strongly enough, but whether that growth can coexist with the UHE
neutrino observations.

This manuscript version records the canonical physics checkpoint reached in
v0.7--v0.8.  A later v0.9 public-data diagnostic uses IceCube's released
energy-dependent differential limits, but it does not alter the meaning of the
v0.8 result documented here.

\section{Literature and context}
\subsection{$\nu_\tau$ regeneration and coupled transport}
The classic treatment of high-energy $\nu_\tau$ propagation through the Earth
was developed by Dutta, Reno and Sarcevic, who emphasized that tau production
and decay regenerate the incident flavor and can substantially reshape hard
spectra \cite{Dutta:2000}.  They subsequently studied secondary neutrinos from
tau decay and showed that the induced $\nu_\mu$ component can matter for
upward-muon observables in sufficiently hard spectra \cite{Dutta:2002}.
Rakshit and Reya formulated and tested iterative solutions of the coupled
$\nu_\tau$--$\tau$ transport equations, including secondary-flavor production
and convergence properties \cite{Rakshit:2006}.  These works provide the
historical benchmarks used to validate the numerical transport basis of the
present project.

The essential physical lesson is that a BSM interaction must be inserted at
the level of the differential transport kernel, not only as a rescaling of a
final event rate.  Changes in total interaction probability, inelasticity,
regeneration, and flavor/topology can compete.  The early stages of this
project therefore separated propagation effects from detector-response
proxies and treated them as distinct diagnostics.

\subsection{Neutrino DIS at ultra-high energy}
At TeV--EeV energies neutrino propagation depends on deep-inelastic scattering
at increasingly small Bjorken $x$.  Modern Standard-Model calculations include
higher-order QCD, modern PDFs, small-$x$ extrapolation and nuclear effects; a
recent CT18-based analysis gives a useful summary of the associated UHE
uncertainties \cite{Xie:2024}.  The present Case-6 calculation uses CT10nlo at
leading order because the immediate goal is a controlled ratio benchmark for
the new operator and a direct connection to the earlier phenomenological
study.  It is not presented as a precision Standard-Model cross-section
prediction.

\subsection{Momentum-growing neutrino--quark interactions}
Pandey, Karmakar and Rakshit analyzed several higher-dimensional neutrino
interactions and compared collider and IceCube constraints \cite{Pandey:2019}.
Their Case 6 contains the derivative dimension-seven structure used here,
\begin{equation}
 {\cal L}_{\rm C6}=C_4
 (\bar\nu_i^c\overleftrightarrow{\partial}_\mu\nu_j)
 (\bar q\gamma^\mu q)+{\rm h.c.},
 \qquad [C_4]=-3.
 \label{eq:C6}
\end{equation}
The benchmark coefficient is
\begin{equation}
 |C_4|=8.6\times10^{-10}\;{\rm GeV}^{-3},
 \label{eq:C4}
\end{equation}
corresponding to the monojet bound quoted in that work.  Here this number is
used as a benchmark rather than as a claim that a contact EFT can be
extrapolated without qualification to every partonic energy.

The distinctive feature is the positive power of the partonic invariant
energy.  This is precisely what was missing from the smoother interactions
examined earlier in the project.  It also creates the central tension of the
present paper: the growth that makes Earth propagation observable does not
stop automatically when the neutrino energy reaches the cosmogenic/UHE
regime.

\subsection{IceCube, KM3NeT and UHE neutrinos as constraints}
IceCube's 12.6-year EHE search constrains neutrino fluxes above the PeV scale
and provides a public true-energy effective area and an energy-dependent
quasi-differential limit \cite{IceCube:2025EHE,IceCube:2025Data}.  The public
release is sufficient for controlled true-energy folds, but not for a new-
physics reconstruction of the collaboration's full likelihood because the
reconstructed-energy/direction response, background templates and nuisance
model are not released in the form required for that task.

KM3NeT's ultra-high-energy event KM3-230213A motivates a complementary
response check because the collaboration released a sky-averaged bright-track
effective area and a public notebook that can be reproduced directly
\cite{KM3NeT:2025Nature,KM3NeT:2025Data}.  The broader ARCA performance study
keeps track and shower channels distinct \cite{KM3NeT:2024ARCA}.  This
separation is crucial here because Case 6 is neutral-current-like: the hard
interaction itself does not create the charged muon selected by the published
bright-track response.

For the UHE stress test we use one public Ehlert cosmogenic-neutrino spectrum
as a fixed spectral shape \cite{Ehlert:2024}.  This is a convenient benchmark
for exposing the observational clash.  It does not change the scientific goal
of the project into a cosmic-ray-composition inference problem.

\section{Momentum-growing Case-6 calculation}
\subsection{Partonic kinematics and matrix element}
For a nucleon at rest, Bjorken $x,y$ give
\begin{equation}
 s\simeq2m_NE_\nu,\qquad \hat s=xs,
 \qquad Q^2=xys=\hat s y.
 \label{eq:kin}
\end{equation}
For the Case-6 normalization used in the v0.7 implementation,
\begin{equation}
 \overline{|{\cal M}_{\rm C6}|^2}
 =64 |C_4|^2\hat s^3 y(1-y),
 \label{eq:M2}
\end{equation}
and the DIS fold is performed explicitly over $x$ and $y$ with the CT10nlo
PDFs.  The factor $y(1-y)$ vanishes at the endpoints, while the $\hat s^3$
growth drives a rapid rise before the EFT-domain restriction is applied.

Two cut prescriptions were retained deliberately:
\begin{align}
 \sqrt{\hat s}<\Lambda &\qquad {\rm(primary)},
 \label{eq:shatcut}\\
 Q<\Lambda &\qquad {\rm(aggressive\ envelope)}.
 \label{eq:Qcut}
\end{align}
The second is weaker because small $y$ can keep $Q$ below $\Lambda$ even when
$\sqrt{\hat s}$ has already exceeded it.  We therefore use
Eq.~\eqref{eq:shatcut} for the primary interpretation.

Defining $c_4=C_4\Lambda^3$, the four benchmark cutoffs give
\begin{equation}
 {c_4\over4\pi}=0.231,\ 0.547,\ 0.994,\ 1.848
 \quad {\rm for}\quad
 \Lambda=1.5,\ 2.0,\ 2.44,\ 3.0\ {\rm TeV}.
 \label{eq:c4pert}
\end{equation}
Thus $2.44$ TeV is already near the nominal $4\pi$ coefficient diagnostic and
$3$ TeV is beyond it.  This is only a perturbativity diagnostic, not a full
partial-wave unitarity analysis.

\subsection{Cross-section result}
The no-cut validation ratios are
\begin{equation}
 {\sigma_{\rm NP}\over\sigma_{\rm SM}}
 =5.34\times10^{-4}\quad(0.1\ {\rm PeV}),
 \qquad
 =1.38\times10^{-2}\quad(1\ {\rm PeV}).
 \label{eq:nocut}
\end{equation}
The conservative $\sqrt{\hat s}<\Lambda$ results are summarized in
Table~\ref{tab:crosssection}.

\begin{table}[t]
\centering
\caption{Case-6 cross-section ratio
$r=\sigma_{\rm NP}/\sigma_{\rm SM}$ with the conservative
$\sqrt{\hat s}<\Lambda$ restriction.}
\label{tab:crosssection}
\begin{tabular}{@{}rrrrrr@{}}
\toprule
$\Lambda$ [TeV] & $c_4/4\pi$ & $r(1\,{\rm PeV})$ &
$r(10\,{\rm PeV})$ & $r(100\,{\rm PeV})$ & $r(300\,{\rm PeV})$\\
\midrule
1.50 & 0.231 & 0.0138 & 0.0909 & 0.148 & 0.177\\
2.00 & 0.547 & 0.0138 & 0.194  & 0.351 & 0.434\\
2.44 & 0.994 & 0.0138 & 0.302  & 0.569 & 0.719\\
3.00 & 1.848 & 0.0138 & 0.404  & 1.03  & 1.33\\
\bottomrule
\end{tabular}
\end{table}

This table establishes the first major physics conclusion.  A genuinely
momentum-growing interaction can evade the qualitative failure of the smooth
benchmarks: it is negligible around a PeV but becomes tens of percent to
order unity by $100$--$300$ PeV.  The earlier smooth-interaction conclusion is
therefore not a robust statement about all BSM $\nu_\tau$ interactions.

\section{The UHE/cosmogenic observational tension}
\subsection{Public-response validation}
The v0.8 stage folded the saved Case-6 curves through the official IceCube EHE
and KM3NeT bright-track public responses.  The KM3NeT calculation reproduces
the collaboration's Ehlert-spectrum benchmark:
\begin{align}
 N_{\rm KM3}^{335\,d}&=0.242754
 \quad {\rm versus}\quad 0.242767\ {\rm(public\ notebook)},\\
 N_{\rm KM3}^{\rm central\ 90\%}&=0.211081
 \quad {\rm versus}\quad 0.211117.
 \label{eq:kmvalidation}
\end{align}
This agreement is an important validation of the interpolation, flux
convention and livetime normalization.

The first IceCube v0.8 fold contained a parser error.  The official response
file is headerless and its area is given in ${\rm m}^2$.  Treating the first
numerical line as column names caused the ${\rm m}^2\to{\rm cm}^2$ conversion
to be skipped, yielding the false value $0.001767$ events.  After the parser
and NumPy read-only-array issue were corrected, the Standard-Model Ehlert fold
became
\begin{equation}
 N_{\rm IC}^{\rm SM}=17.674595
 \qquad(12.6\ {
m yr}).
 \label{eq:icsm}
\end{equation}
The failed value and its correction are retained in the project history.

\subsection{Corrected v0.8 event-proxy result}
Using the conservative cross-section enhancement as an inclusive IceCube
response proxy gives Table~\ref{tab:uheproxy}.

\begin{table}[t]
\centering
\caption{Corrected v0.8 IceCube EHE response proxy for the unscaled Ehlert
benchmark.  These are not collaboration-level BSM detector predictions.}
\label{tab:uheproxy}
\begin{tabular}{@{}rrrr@{}}
\toprule
$\Lambda$ [TeV] & $N_{\rm IC}^{\rm SM}$ &
$N_{\rm IC}^{\rm SM+NP,proxy}$ & fractional increase\\
\midrule
1.50 & 17.67 & 20.73 & 0.173\\
2.00 & 17.67 & 25.05 & 0.417\\
2.44 & 17.67 & 31.15 & 0.763\\
3.00 & 17.67 & 42.54 & 1.407\\
\bottomrule
\end{tabular}
\end{table}

The aggressive $Q<\Lambda$ envelope would instead give approximately
$37.73$, $61.93$, $92.85$ and $145.9$ events and is therefore even more
strongly in tension with the UHE response.  The important point is not the
precise proxy count but the direction of the effect: the same momentum growth
that makes the Earth-propagation signal interesting continues into the UHE
regime where the experimental constraint is severe.

\subsection{Why a common-flux normalization largely cancels the enhancement}
Let $I$ and $K$ denote the IceCube and KM3NeT yields per unit common flux
normalization, including whatever response model is being tested.  Then
\begin{equation}
 N_{\rm IC}=\kappa I,\qquad
 N_{\rm KM3}=\kappa K.
 \label{eq:commonflux}
\end{equation}
Any upper bound of the form $\kappa\le U/I$ implies
\begin{equation}
 N_{\rm KM3}\le U{K\over I}.
 \label{eq:ratioidentity}
\end{equation}
If the same momentum-growing interaction raises $I$ and $K$ by comparable
factors, the ratio $K/I$ changes much less than either absolute yield.

At the v0.8 checkpoint we used only a deliberately crude zero-count Poisson
screen,
\begin{equation}
 \kappa_{90}^{\rm screen}={2.302585\over N_{\rm IC}^{\rm proxy}},
 \label{eq:crudescreen}
\end{equation}
not the published IceCube likelihood.  Under that common normalization, the
KM3NeT \emph{NC-equivalent} response proxy becomes
\begin{equation}
 N_{\rm KM3}^{\rm NC\mbox{-}equiv}
 \simeq 0.0315,\ 0.0314,\ 0.0312,\ 0.0310
 \label{eq:kmproxy}
\end{equation}
for $\Lambda=1.5,2.0,2.44,3.0$ TeV.  The near-independence of the cutoff is
therefore not accidental; it follows from the common-flux structure.

\subsection{Topology: the Case-6 operator is not a direct bright-track channel}
Equation~\eqref{eq:C6} has a neutral-current-like hard final state: neutrino
sector plus hadrons.  It does not directly produce the charged muon selected
by the public KM3NeT bright-track response.  Consequently
$A_{\rm track}^{\rm SM}(1+r)$ is only an acceptance proxy and must not be
called a physical track prediction.

At the crude v0.8 screening scales, the unchanged direct bright-track
expectation was approximately
\begin{equation}
 0.027,\ 0.022,\ 0.018,\ 0.013
 \label{eq:directtrack}
\end{equation}
for the same four cutoffs.  A genuine Case-6 track effect would have to arise
indirectly through modified propagation, flavor structure, or secondary
neutrinos and therefore requires the coupled transport calculation rather
than a multiplicative track-area rescaling.

\section{What the v0.8 checkpoint establishes}
The physics conclusion at this stage is deliberately narrow:
\begin{quote}
A genuinely momentum-growing interaction can produce the large high-energy
cross-section modification needed to make $\nu_\tau$ propagation interesting,
but the same growth creates a serious UHE/cosmogenic observational tension
when folded through the public IceCube and KM3NeT responses.
\end{quote}
This is not a no-go theorem.  It is the identification of the central tension
that any viable momentum-enhanced $\nu_\tau$ propagation model must resolve.

Several limitations prevent a stronger event-level statement.  First,
Eq.~\eqref{eq:crudescreen} ignores the three PeV-scale IceCube events and the
energy/direction information used by the collaboration.  Second, the released
IceCube effective area already contains Standard-Model propagation and
selection effects; multiplying it by an inclusive cross-section ratio does
not reproduce a detector simulation with altered NC visible energy,
attenuation or regeneration.  Third, a tau-specific operator cannot simply
rescale the full all-flavor IceCube area.  Fourth, the KM3NeT bright-track area
is not the acceptance for the Case-6 NC final state.

The next step fixed at this checkpoint is therefore to replace the crude
$2.3026$-event screen with the published energy-dependent IceCube EHE
information as faithfully as the public data allow, while refusing to call a
true-energy effective-area fold the collaboration's exact binned likelihood.
That calculation became the subsequent v0.9 public-data diagnostic.

\section{Reproducibility and version history}
The historical transport, smooth-interaction scans, Case-6 PDF integration,
and UHE response stages are versioned in the repository.  The failed IceCube
unit conversion and its correction are retained explicitly rather than being
silently removed.  The canonical numerical inputs for this manuscript are the
saved v0.7/v0.8 Case-6 cross-section table and the public IceCube/KM3NeT
responses.  Recomputing the LHAPDF grid is unnecessary for the v0.8 event fold.

The project log must preserve the distinction between three objects:
(i) the physical Earth-propagation calculation, (ii) detector-response proxies,
and (iii) collaboration likelihoods.  No statement based on (ii) is promoted
to (iii).  Likewise, the use of a cosmogenic spectral benchmark as a UHE
stress test is not a change of project scope.

\section{Conclusions}
The momentum-growing Case-6 operator changes the trajectory of the
$\nu_\tau$-propagation project.  Smooth benchmark interactions had suggested
that observable distortions might remain too weak after nuisance profiling.
The full PDF-level Case-6 calculation shows that this is not a general
conclusion: within the conservative EFT-domain cut, the new contribution can
reach tens of percent to order unity at $100$--$300$ PeV while remaining small
near a PeV.

The price is a new and more stringent problem.  The same momentum growth
extends into the UHE regime probed by IceCube, and the corrected public-response
fold exposes a strong observational tension.  At the v0.8 level this is only a
screening result, because neither the full BSM detector response nor the exact
IceCube reconstructed-bin likelihood is publicly reproducible from the
released tables.  Nevertheless, it identifies the correct next physics task:
retain the momentum-enhanced $\nu_\tau$ propagation effect while satisfying
the UHE observational constraint in a model whose differential transport and
detector topology are treated consistently.

\section*{AI-assisted drafting disclosure}
This working manuscript and its reproducibility notes were prepared with AI
assistance under the author's direction.  The author should independently
verify every numerical and bibliographic statement before submission.

\bibliographystyle{../v0.9/JHEP}
\bibliography{../v0.9/references}
\end{document}
